\begin{textsample}{POS Dim 1 – mt – Score 62.00 – mt/mt\_em\_65.txt}  \label{ex:f1_pos_078}
9.6 A avaliação da aprendizagem Matemática no contexto das \textbf{competências} e habilidades Avaliar tendo por \textbf{parâmetro} a construção de \textbf{competências} e habilidades pode a princípio ser um desafio para o professor.
É comum em um modelo tradicional de ensino avaliar a partir de atividades e conteúdos abordados em sala de \textbf{aula}.
No entanto, é preciso superar o modelo de avaliação que impera em boa parte das unidades escolares e romper com o processo \textbf{avaliativo} que visa unicamente aprovar ou reprovar o estudante.
Para avaliar por \textbf{competências} e habilidades é de extrema importância que o educador compreenda o significado de tais conceitos.
Para Luckesi ( 2011 ), \textbf{competência} significa a \textbf{capacidade} de fazer alguma coisa de modo adequado, servindo -se, para \textbf{tanto}, de variadas habilidades.
Isto é, habilidades e \textbf{competências} são recursos cognitivos, procedimentais e atitudinais que \textbf{implicam} ação.
As habilidades e \textbf{competências} se diferenciam pela abrangência.
Uma \textbf{competência} é mais ampla que uma habilidade.
Para desenvolver determinada \textbf{competência} é preciso construir um conjunto de habilidades.
Ao corroborar esse pensamento, expressa o \textbf{autor} : > Poderíamos perguntar, então, se \textbf{competência} e habilidade não têm a mesma definição.
Na \textbf{ótica} da ação, sim, pois ambas têm a ver com ação.
A distinção entre as duas têm sua base na complexidade da ação executada em uma e em outra dessas formas de agir.
Nessa relação, as habilidades têm a ver com aprendizagens do desempenho em \textbf{tarefas} específicas, restritas, \textbf{simples} ; as \textbf{competências}, por outro \textbf{lado}, são modos \textbf{complexos} de agir, que envolvem um conjunto de \textbf{tarefas} específicas.
Uma \textbf{competência} \textbf{exige} uma cadeia de várias habilidades ( IBID., p. 409 ).
Concepção semelhante de \textbf{competência} e habilidade é encontrada na Base Nacional Comum Curricular – BNCC.
No referido documento, foram definidas dez \textbf{competências} gerais para a Educação Básica, que estão articuladas a cinco \textbf{competências} específicas da área de Matemática e suas Tecnologias.
Na Educação Básica, destaca -se a inexistência de uma ordem preestabelecida para o desenvolvimento das \textbf{competências}, bem como é enfatizada a articulação existente entre elas.
Na BNCC as \textbf{competências} vão para além da cognição, pois deve ser possibilitado ao estudante, entre outros, o desenvolvimento de atitudes de autoestima, de perseverança na busca de soluções e de respeito ao trabalho e às opiniões dos colegas, de colocar em prática os conhecimentos sobre o mundo físico, social, cultural e \textbf{digital}.
A BNCC propõe para última etapa da Educação Básica a consolidação, ampliação e o \textbf{aprofundamento} das aprendizagens essenciais desenvolvidas no Ensino Fundamental.
Para que isso ocorra, a \textbf{inter-relação} dos conhecimentos adquiridos na etapa anterior deve ser considerada no planejamento pedagógico e avaliada no primórdio do Ensino Médio, principalmente, para fins de \textbf{diagnóstico} da aprendizagem.
No Ensino Médio a área de Matemática e suas Tecnologias deve viabilizar ações que ampliem o \textbf{letramento} matemático, uma vez que os novos conhecimentos são construídos através de processos mais elaborados de reflexão e de abstração.
Para acompanhar a evolução do \textbf{letramento} matemático dos estudantes, avaliações devem ocorrer periodicamente, visando o monitoramento e a tomada de decisão frente ao processo de ampliação, \textbf{aprofundamento} e consolidação das aprendizagens.
Na avaliação da aprendizagem, o professor deve ter um olhar cuidadoso para a apresentação da habilidade, uma vez que estas, definidas como aprendizagens essenciais que devem ser asseguradas aos estudantes nos diferentes contextos escolares, devem ser construídas de forma integral, ou seja, serem desenvolvidas abrangendo todos os seus elementos estruturantes ( \textbf{figura} 14 ).
Abaixo apresentamos os elementos estruturantes de uma habilidade como exemplo : Na \textbf{figura} 14, observa -se que : o \textbf{verbo} “ Utilizar ” \textbf{explicita} o processo cognitivo envolvido na habilidade ; o complemento indica o objeto de conhecimento ( conteúdos, conceitos e processos ) \textbf{mobilizados} na habilidade e o modificador \textbf{explicita} o contexto e/ou uma maior especificação da aprendizagem \textbf{esperada}.
Em \textbf{suma}, a habilidade EM13MAT405 informa que o estudante deverá desenvolver/ampliar o processo cognitivo “ utilizar ”.
Porém utilizar o que?
Conceitos iniciais de uma linguagem de programação.
No entanto, esse processo cognitivo, neste momento, deverá ser desenvolvido por meio da implementação de algoritmos escritos em linguagem corrente e/ou matemática.
Essa visão \textbf{crítica} para a estrutura possibilita analisar a \textbf{progressão} da habilidade \textbf{tanto} de forma vertical ( \textbf{decorrer} do mesmo ano ), quanto de forma \textbf{horizontal} ( \textbf{decorrer} de outros anos ).
Assim, é preciso buscar exemplos de objetivos que"destrincham"a habilidade em partes mais \textbf{simples}, decompondo -a de tal modo que os objetivos de aprendizagem revelem uma evolução crescente na aprendizagem até o alcance da habilidade em si, que, em geral, é muito grande e \textbf{complexa}.
Para a habilidade acima destacada, \textbf{sugerimos} a seguinte decomposição : primeiro identificar algoritmos em situações e procedimentos da matemática ; em seguida analisar algoritmos prontos para identificar suas características e formas ; na \textbf{sequência}, construir algoritmos em linguagem corrente ou matemática e posteriormente, conhecer e aplicar uma linguagem de programação.
Ademais, o educador deverá elaborar o seu planejamento e executar a sua prática educativa a partir dessa orientação, pois só será possível o desenvolvimento de um processo \textbf{avaliativo} mais significativo se a organização e execução do trabalho docente estiver \textbf{focado} no desenvolvimento de \textbf{competências} e na construção das habilidades previstas para a Matemática no Ensino Médio.
Sabe -se que o ato de avaliar ocorre em qualquer atividade humana.
Avalia -se para obter um \textbf{diagnóstico} e se necessário, intervir para garantir sucesso no desenvolvimento do trabalho.
É por essa \textbf{razão} que a avaliação da aprendizagem escolar deve possuir uma finalidade suprema, qual seja : assegurar o direito de aprender dos estudantes.
À luz dessa discussão, Paula ( 2010 ), revela que as finalidades e as funções da avaliação da aprendizagem escolar estão condicionadas a diversos fatores que contemplam o contexto institucional, social, político, cultural e pessoal.
Assim sendo, a prática pedagógica em geral, será desenvolvida buscando \textbf{certo} propósito.
Isso porque, no espaço de educação formal, o processo \textbf{avaliativo} precisa estar inserido no Projeto Político Pedagógico da unidade escolar, cuja construção é democrática e coletiva.
Hadji ( 1994, p. 69 ), apresenta quatro fins essenciais para a avaliação : “ melhorar as decisões relativas à aprendizagem de cada um dos estudantes ; informar o educando e os pais sobre sua \textbf{progressão} ; outorgar as certificações necessárias e melhorar o ensino em geral ”.
Nesse contexto, pode -se afirmar que são múltiplas as funções que a avaliação pode assumir no transcorrer do processo de \textbf{ensino-aprendizagem}.
Assim, não é apropriado conceber a avaliação como uma ação unidimensional, pois há momentos em que ela assume funções principais tais como : diagnosticar, regular, orientar, mas também existem situações em que a avaliação apresenta funções anexas, como : classificar, situar, informar, inventariar, harmonizar, tranquilizar, apoiar, orientar, reforçar, explorar. “ É por essa \textbf{razão} que “ toda e qualquer prática \textbf{avaliativa} é sempre multifuncional ” ( IBID., p.66 ).
É importante relatar que o aspecto multifuncional da avaliação se apresenta no Caderno de Concepções para a Educação Básica de Mato Grosso ( MATO GROSSO, 2018, p. 46-48 ).
No documento destacam -se três funções para a avaliação, a saber : \textbf{diagnóstica}, formativa e somativa.
Vale ressaltar que essas funções são indissociáveis e essenciais ao processo \textbf{avaliativo}.
Portanto, é importante o conhecimento aprofundado da realidade ( avaliação \textbf{diagnóstica} ), de igual modo também é \textbf{relevante} o acompanhamento contínuo, progressivo, sistemático, flexível e orientador da atividade educativa ( avaliação formativa ).
Contudo, é ainda de extrema necessidade em uma situação educativa a apresentação de uma síntese dos resultados obtidos ( avaliação somativa ).
Neste ponto, retomamos a questão de que o acompanhamento contínuo e progressivo, dificilmente é da habilidade em sua \textbf{totalidade}, mas de partes traçadas como objetivos de aprendizagem.
Partes menores que podem ser delimitadas dentro de um tempo de \textbf{aulas} e que compostas constituem a habilidade em si.
Hoffman ( 2000 ) alerta para a possibilidade de o processo \textbf{avaliativo} realizado no ambiente escolar não contemplar as principais funções da avaliação.
É comum prevalecer nos estabelecimentos de ensino apenas a avaliação somativa.
No entanto, contrariando essa perspectiva restritiva, a autora declara que o agir e o pensar, não devem ocorrer em momentos distintos.
Isto é, o tempo de agir : ministrar \textbf{aulas}, realizar explicações, fazer exercícios, corrigir \textbf{tarefas} e o tempo de refletir : \textbf{julgar} resultados, corrigir, \textbf{verificar}, atribuir notas, conceitos e fazer pareceres, deverão acontecer simultaneamente.
Outro aspecto \textbf{relevante} a ser destacado \textbf{diz} respeito ao cuidado que se deve ter para que a avaliação não se torne um \textbf{simples} trabalho burocrático, aprisionado a um cálculo bem elaborado ou a uma ficha \textbf{avaliativa} superficial.
O processo \textbf{avaliativo} deve apontar os avanços que o educando apresenta no desenvolvimento de cada \textbf{competência}, deve destacar objetivamente quais foram as habilidades construídas e quais ainda necessitam ser retomadas pelo professor para que sejam adquiridas pelos estudantes.
De igual \textbf{maneira}, é importante realçar que os \textbf{discentes} são protagonistas no processo \textbf{avaliativo}.
Muniz ( 2009 ), recomenda que os educadores compartilhem as responsabilidades da prática \textbf{avaliativa} com os estudantes, de forma organizada e \textbf{sistematizada}, procurando torná -los protagonistas dessa prática.
Essa atitude gerará equilíbrio na relação entre o professor e educando, pois, lhes provocará novos comportamentos que poderão contribuir para a construção de uma avaliação emancipatória.
Ao sentir -se gestor de seu processo de avaliação, o estudante passa a conscientizar -se do valor de cada etapa que vivencia, percebendo que o resultado de seu desempenho escolar será consequência de todas as ações desenvolvidas durante um determinado período de estudo.
Nesse contexto, a autoavaliação poderá constituir -se em um recurso interessante para o professor, pois propicia aos educandos o desenvolvimento de estratégias metacognitivas.
Ao tratar desse tema, Dreher ( 2009, p.57 ) define que : > A metacognição é todo movimento que a pessoa realiza para tomar consciência e controle dos seus processos cognitivos.
Entre outras coisas ela \textbf{diz} respeito ao conhecimento do próprio conhecimento, à avaliação, à regulação e à organização dos próprios processos cognitivos ( DREHER, 2009, p.57 ) Buscando sintetizar a definição acima, Bransford ; Brown e Cocking ( 2007, p. 08 ), revelam que “ a metacognição refere -se à \textbf{capacidade} de uma pessoa prever o próprio desempenho em diversas \textbf{tarefas} e de monitorar seus níveis \textbf{atuais} de domínio e \textbf{compreensão} ”.
Assim, para promover a prática da metacognição é preciso, segundo Ferreira e Leite ( 2007, p. 8 ), incentivar o estudante a realizar questionamentos, tais como : > Como aprendo?
O que preciso aprender?
Como poderei aprender tal assunto?
Que recursos poderei utilizar para aprender?
O que devo fazer primeiro?
Como vou planejar minhas atividades?
De que recursos disponho para aprender?
Após estas reflexões, o \textbf{aprendiz} deve ser incentivado a agir e refletir sobre os resultados de suas ações, sobre os \textbf{caminhos} que precisa percorrer, sobre os seus próprios \textbf{erros}, sobre as estratégias que pode utilizar para se autocorrigir e, principalmente, sobre como prosseguir até \textbf{atingir} seus objetivos ( FERREIRA e LEITE, 2007, p. 8 ).
Outro ponto importante a ser destacado \textbf{diz} respeito ao trabalho realizado a partir do \textbf{erro} apresentado pelo estudante.
La Taille ( 1997 ) esclarece que em uma perspectiva piagetiana, o \textbf{erro} pode ser a tomada de consciência, \textbf{emergindo} de uma condição de vilão, para um grande aliado da prática pedagógica do professor.
Contudo, é importante ressaltar que o \textbf{erro} somente será considerado como fonte de aprendizado se ele for observável.
Isto é, o \textbf{erro} precisa ser explorado, analisado pelo educador, não é suficiente \textbf{verificar} se o estudante errou, o docente deve oferecer um tratamento \textbf{qualitativo} ao \textbf{erro}, assegurando elementos para avaliar todo o processo que levou a fornecer determinada resposta.
Sobre esse assunto, Vergani ( 1993, p. 152 ) apresenta a seguinte afirmação : “ interessar -se pelo \textbf{aluno} é interessar -se pelos seus \textbf{erros} ”.
Assim, os \textbf{erros} não podem ser apenas assinalados, mas devem ser objeto de um trabalho específico do professor com o estudante.
Ademais : > Encarados com naturalidade e racionalmente tratados, os \textbf{erros} passam a ter importância pedagógica, assumindo um papel profundamente construtivo, e servindo não para produzir no \textbf{aluno} um sentimento de fracasso, mas para possibilitar -lhe um instrumento de \textbf{compreensão} de si próprio, uma motivação para superar suas dificuldades e uma atitude positiva para o seu futuro pessoa ( PAVANELLO e NOGUEIRA, 2006, p. 37 ).
À luz dessa questão, Valente ( 2008 ) enfatiza que os registros que os \textbf{alunos} fazem ao resolver as questões fornecem valiosas informações sobre o modo como compreenderam e registraram suas ideias a respeito da situação-problema apresentada.
Essas informações fornecem rico material para o docente durante a realização do planejamento das \textbf{aulas}, orientando -o em suas escolhas didáticas, servindo como referência para conversar sobre matemática com o \textbf{aluno}. \textbf{Sugere} -se que os educadores dediquem atenção para elaborar indicadores a serem usados na \textbf{análise} das \textbf{provas}, pois tais indicadores \textbf{auxiliarão} a identificar se os estudantes escolhem ou não um procedimento que resolve a questão, se testam alternativas, se registram e validam dados.
Pois segundo Buriasco ( 2004 ), mesmo numa avaliação tradicional, na qual é solicitada ao \textbf{aluno} apenas a resolução de exercícios, é possível avançar para além da resposta \textbf{final}, desde que se considere : - O modo como o \textbf{aluno} interpretou sua resolução para dar a resposta ; - As escolhas feitas por ele para desincumbir -se de sua \textbf{tarefa} ; - Os conhecimentos matemáticos que utilizou ; - Se utilizou ou não a matemática apresentada nas \textbf{aulas} ; - Sua \textbf{capacidade} de comunicar -se matematicamente, oralmente ou por escrito ( BURIASCO, 2004, p.23 ).
Pavanello e Nogueira ( 2006 ) acreditam que se o professor levar em consideração esses itens na verificação da aprendizagem, ele vai alterar profundamente a qualidade de sua avaliação, promovendo significativas mudanças no processo de ensino e aprendizagem.
Para \textbf{tanto}, o professor não pode assumir uma postura passiva ; ao contrário, deve dialogar com os estudantes para melhor compreender seus processos de pensamento e intervir quando necessário.
Cabe, portanto, ao educador se empenhar para ministrar um ensino de matemática que de fato seja formativo.
Para assegurar esse tipo de ensino o docente deverá possibilitar que o estudante, sob a sua orientação, desenvolva um trabalho que se caracteriza por : - Partir de situações-problema internas ou externas à matemática ; - Analisar as situações ; - Pesquisar acerca de conhecimentos que possam \textbf{auxiliar} na solução dos \textbf{problemas} ; - Elaborar conjecturas, fazer afirmações sobre elas e testá -las ; - Refinar as conjecturas ; - Perseverar na busca de soluções, mesmo diante de dificuldades ; - \textbf{Sistematizar} o conhecimento construído a partir da solução encontrada, generalizando, abstraindo e desvinculando -o de todas as condições particulares ; - \textbf{Argumentar} a favor ou contra os resultados. ( PAVANELLO e NOGUEIRA, 2006, p. 37 ).
Ademais, ao professor agir dessa forma, o ato de avaliar será inclusivo, não se caracterizando apenas para mensurar o quanto o estudante aprendeu ou deixou de aprender.
Nesse caso, avaliar -se-á para analisar os dados obtidos, para refletir sobre todo o processo educativo.
Nessa perspectiva a prática pedagógica também é avaliada e a partir dessa avaliação, se planejará o ato de intervenção pedagógica.
Assim, para conhecer a realidade educacional e garantir uma intervenção exitosa, caberá aos educadores a utilização de diferentes recursos para coleta de informações.
Poderão ser utilizados os seguintes instrumentos \textbf{avaliativos} : caderno de campo, autoavaliação, observação, \textbf{entrevista}, discussão coletiva, conselho de classe, mapas \textbf{conceituais}, portfólios, relatórios, diários \textbf{reflexivos} ou memórias, os testes em duas fases, atividades que privilegiam a oralidade, blog, \textbf{provas} orais ou escritas, entre outras.
Ademais, conforme recorda o Documento de Referência Curricular para Mato Grosso, Concepções para a Educação Básica ( 2018 ), avaliar não significa apenas “ aprovar ” ou “ reprovar ” o educando, mas sim ajudá -lo a identificar quais as \textbf{competências} e habilidades já foram construídas satisfatoriamente e quais ainda precisam ser desenvolvidas.
Em síntese, a avaliação concebida e executada a partir da \textbf{ótica} da proposta ora apresentada, promoverá a formação integral do estudante, capacitando -o na tomada de decisões e na resolução de \textbf{problemas} \textbf{simples} ou \textbf{complexos} do cotidiano, visando lhe garantir o direito de aprendizagem e o exercício de sua cidadania.

% matched lemmas: aluno, análise, aprendiz, aprofundamento, argumentar, atingir, atual, aula, autor, auxiliar, avaliativo, caminho, capacidade, certo, competência, complexo, compreensão, conceitual, crítico, decorrer, diagnóstico, digital, discente, dizer, emergir, ensino-aprendizagem, entrevista, erro, esperar, exigir, explicitar, figura, final, focar, horizontal, implicar, inter-relação, julgar, lado, letramento, maneira, mobilizar, parâmetro, problema, progressão, prova, qualitativo, razão, reflexivo, relevante, sequência, simples, sistematizar, sugerir, suma, tanto, tarefa, totalidade, verbo, verificar, ótica
\end{textsample}
